\chapter{Discussion}
\label{chap:discussion}

\section{The Role of Memory Representation}
Across the diagnostic and full runs, this small generator follows short code fragments more reliably than accumulated textual lessons. The text and code systems differ in more than representation, so this is a cross-version observation rather than a controlled causal comparison. It is consistent with Metis's finding that code memory can complement text memory~\cite{dai2026metis}, while adding a setting in which the format of the memory may determine whether a 7B generator can use it at all. The form of code also matters: complete files were sometimes copied as finished answers, whereas fragments prevented undeclared imports less reliably.

This pattern differs from reports for stronger models, which transfer text memory more reliably than code~\cite{dai2026metis,kim2026mtl}. The mechanisms are not identical because Metis executes code whereas this method displays code to the generator. Taken together, the evidence suggests that an appropriate memory representation depends on the capability of the model consuming it.

Part of the observed gain is protection against collapse rather than more capable application logic. In an early phase, short prompts led the model to return the default scaffold in five of eight tasks, whereas longer, specific prompts did not. Both effects can improve system performance, but they should be reported separately.

\section{What Is Learned and at What Cost}
The generator itself does not change. What changes is the code shown to it and the conditions under which it is shown. The resulting behaviour is closer to accumulating a small set of preventive code examples than to acquiring programming ability. It is also a property of the entire system, including a memory model that is larger than the generator. The memory remains small across 101 tasks; the study does not test what happens over thousands of tasks, when more problem types compete for three slots.

The system also adds computation. A Long Memory task takes approximately 12 minutes, compared with approximately 6 minutes for No Memory, largely because the 72B model constructs checklists, judges applicability, and reflects after each task. If the fixed-prefix ablation recovers most of the gain, a deployed system could use the learned fragments without invoking the 72B model at generation time, at the cost of losing online adaptation.

\section{Limitations and Threats to Validity}
\label{sec:threats}
\begin{itemize}
  \item \textbf{Single runs.} WebGen-Bench itself reports single runs without error bars. No Memory is repeated under several orders to estimate run-to-run noise, and Long Memory is repeated under further orders using paired comparisons. Within any order, however, Long Memory remains a single run. Two runs of the same Long Memory version can differ by approximately five points, which is more variation than observed for No Memory. The task-level intervals therefore understate full uncertainty. Orders share the same 101 tasks and are not independent samples. \todo{Status: a second manifest run of the final method is queued (about 10-19); report the difference between the two runs.}
  \item \textbf{Gain beyond the checklist and trade-offs between forms.} The final method's paired gain over checklist only is $+3.7$ [$-1.8$, 9.1], so an effect beyond requirement decomposition is not established. Whole-file examples also prevent undeclared imports more reliably in the original order, so neither code representation dominates everywhere.
  \item \textbf{One generator and one scaffold.} All results concern WebGen-LM-7B, Bolt.diy, and a Vue template. Larger generators may follow textual advice better and gain less from code examples.
  \item \textbf{Code pool as a learned prefix.} The general triggers make the code pool partly resemble a fixed prompt prefix. The fixed-prefix ablation (\Cref{sec:ablation}) is intended to measure how much of the effect this accounts for.
  \item \textbf{Rare problems are learned late.} A safeguard can appear only after the relevant problem has occurred. Memory may therefore help frequent problems more than rare ones, and several problem types can compete for three slots (\Cref{sec:fix}).
  \item \textbf{Evaluators.} Appearance uses Qwen2.5-VL-32B rather than GPT-4o and is therefore compared only across the present arms. \gls{gui} testers can exhaust their interaction budget or misjudge pages~\cite{zheng2023judge}. Raw transcripts are retained, but the judgements were not independently reviewed by a human.
  \item \textbf{Iterative development on the same benchmark.} This is the most serious limitation. Every version was designed after failures were observed on the same 101 tasks used for evaluation. The effect is visible: whole-file code memory, developed on the manifest order, scores 40.0\% there but 33.1\% and 32.6\% in the two shuffled orders, while No Memory stays at 27.9--28.6\%. Empty memories and freezing the final version before all runs except seed 11 reduce this risk but do not remove it. A held-out evaluation is required.
  \item \textbf{Memory model stronger than generator.} The 72B memory model exceeds the 7B generator in size, so some gains may reflect the larger model's organisation of the task. The checklist only restates the requirement, every fragment is copied verbatim from an earlier site generated by the 7B model, and the memory model never sees the current task's evaluation. The results therefore concern what a small generator gains when a stronger model organises its prior experience, rather than self-improvement in isolation. \todo{Status: fixed-prefix run queued after seed 47. State how much of the gain remains without accumulated memory and without the 72B model at generation time.}
\end{itemize}

\section{Ethical Considerations and Impact}
\paragraph{Data and privacy.} The benchmark requirements are synthetic and contain no personal data. All models run on university hardware, and no task data are sent to external services.
\paragraph{Environmental cost.} A full run occupies two H200 GPUs for about one day (\Cref{sec:design}). To limit unnecessary computation, changes were first checked on three to five tasks and with static code analysis before full evaluation.
\paragraph{Self-improving agents and society.} A memory system that turns outcomes into later behaviour can preserve mistakes as readily as useful practices. This design retains provenance for every memory entry so that injected material can be traced, but systems that store and reuse code in deployment would still require security review. Small local models can reduce website-development costs for smaller organisations, while potentially displacing routine front-end work. Their output remains subject to human review.
\paragraph{Use of AI tools.} Claude and ChatGPT were used to assist with programming the experimental pipeline and analysis scripts, organising experiment logs, and giving suggestions in critical reviews of drafts; Gemini and Microsoft Copilot were also used for review comments. All text was written by the author, who decided which suggestions to adopt. Research decisions and the verification of all results are the author's responsibility; every number was computed by the scripts in the public repository and checked by the author. \todo{Author: this statement holds only after the final text has been rewritten in your own words; check it against what was actually done before submission.}

\paragraph{Research integrity.} We report negative results (text memory's decline, the scaffold failure of whole-file examples) and invalidated our own earlier results when a delivery bug was found.

\section{Project Management}
The project was organised in versions. The cluster ran one frozen version at a time; new ideas were tested locally on small samples and with static code checks, and a version was replaced only when the new one was shown to be better. Decisions and their evidence were recorded continuously in a decision log, which made it possible to reconstruct why each component exists.
